In modern distributed systems, data exchange between servers must often balance two opposing objectives:
\emph{enabling communication} within a trusted domain and \emph{preventing communication} by potential adversaries.
When information flows through a network, it is desirable to establish a reliable path between authorized servers
while simultaneously enforcing structural constraints that deny any alternative communication routes accessible to an opponent.
This dual requirement leads naturally to the study of network configurations that combine
\emph{communication} for legitimate participants with \emph{control} mechanisms that isolate or restrict hostile ones.
In this context, one seeks structures that ensure communication only within a specified subnetwork while suppressing all competing connections.
In this paper, we formalize this concept through the notion of a \emph{cut-path},
a structure that simultaneously contains a connecting path and a separating cut,
and we study the computational and structural properties of the corresponding \MinCutPath{} problem.

